% ===============================
% Worksheet / Manual Template
% ===============================
\documentclass[12pt, a4paper]{article}

% ===============================
% Metadata
% ===============================
\newcommand*{\CourseCode}{MAT321}
\newcommand*{\CourseTitle}{Real Analysis II}
\newcommand*{\Chapter}{Metric in $\mathbb{R}^n$, $\mathbb{C}^n$ and Function Spaces}
\newcommand*{\Topics}{Metric, pseudo-metric and Examples}
\newcommand*{\DocDate}{February 01, 2026}

\include{header}

% ==========================================
% Document
% ==========================================
\begin{document}
\FrontPage
\Large

\begin{heading}
    Motivation of Metric
\end{heading}

\clearpage

\begin{heading}
    Metric Space
\end{heading}

\vspace{10pt}

\begin{definition}[Metric Space]
    A \textbf{metric space} is a pair $(X,d)$ where $X\neq\varnothing$ and
    $d:X\times X\to\mathbb{R}$ satisfies, for all $x,y,z\in X$:
    \begin{enumerate}[label=(M\arabic*)]
    \item $d(x,y)\ge 0$;
    \item $d(x,y)=0 \iff x=y$;
    \item $d(x,y)=d(y,x)$;
    \item $d(x,z)\le d(x,y)+d(y,z)$.
    \end{enumerate}
    Here $d$ is called a \textbf{metric} on $X$.
\end{definition}

\clearpage

\begin{definition}[Pseudo-Metric Space]
    A \textbf{pseudo-metric space} is a pair $(X,d)$ where $X\neq\varnothing$ and
    $d:X\times X\to\mathbb{R}$ satisfies, for all $x,y,z\in X$:
    \begin{enumerate}[label=(PM\arabic*)]
    \item $d(x,y)\ge 0$;
    \item $d(x,y)=d(y,x)$;
    \item $d(x,z)\le d(x,y)+d(y,z)$.
    \end{enumerate}
    Here $d$ is called a \textbf{pseudo-metric} on $X$.
\end{definition}

\clearpage

\begin{heading}
    Metric on Real Line
\end{heading}

\vspace{10pt}

\begin{theorem}[Euclidean Metric on Real Line]
    Show that $d:\mathbb{R}\times\mathbb{R}\to\mathbb{R}$ defined as
    \[
        d(x,y)=|x-y|
    \]
    is a metric on $\mathbb{R}$.
\end{theorem}

\clearpage

\begin{theorem}[Another Metric on Real Line]
    Show that $d:\mathbb{R}\times\mathbb{R}\to\mathbb{R}$ defined as
    \[
        d(x,y)=3|x-y|
    \]
    is a metric on $\mathbb{R}$.
\end{theorem}

\clearpage

\begin{problem}
    Show that $d:\mathbb{R}\times\mathbb{R}\to\mathbb{R}$ defined as
    \[
        d(x,y)=|x-2y|
    \]
    is a not metric on $\mathbb{R}$.
\end{problem}

\clearpage

\begin{problem}
    Show that $d:\mathbb{R}\times\mathbb{R}\to\mathbb{R}$ defined as
    \[
        d(x,y)=(x-y)^2
    \]
    is a not metric on $\mathbb{R}$.
\end{problem}

\clearpage

\begin{theorem}
    Show that $d:\mathbb{R}\times\mathbb{R}\to\mathbb{R}$ defined as
    \[
        d(x,y)=\sqrt{|x-y|}
    \]
    is a metric on $\mathbb{R}$.
\end{theorem}

\clearpage

\begin{heading}
    Metric on $\mathbb{R}^2$
\end{heading}

\vspace{10pt}

\begin{theorem}[Standard Euclidean Metric on $\mathbb{R}^2$]
    The metric space $\mathbb{R}^2$, called the \emph{Euclidean plane}, is obtained by
    taking the set of ordered pairs of real numbers,
    \[
    x=(\xi_1,\xi_2), \qquad y=(\eta_1,\eta_2),
    \]
    and defining the \emph{Euclidean metric} by
    \[
    d(x,y)=\sqrt{(\xi_1-\eta_1)^2+(\xi_2-\eta_2)^2}
    \]
    Show that $d$ is a metric on $\mathbb{R}^2$.
\end{theorem}

\clearpage

\begin{theorem}[Taxicab Metric on $\mathbb{R}^2$]
    The metric space $\mathbb{R}^2$, called the \emph{Euclidean plane}, is obtained by
    taking the set of ordered pairs of real numbers,
    \[
    x=(\xi_1,\xi_2), \qquad y=(\eta_1,\eta_2),
    \]
    and defining the \emph{Euclidean metric} by
    \[
    d(x,y)=|\xi_1-\eta_1|+|\xi_2-\eta_2|
    \]
    Show that $d$ is a metric on $\mathbb{R}^2$.
\end{theorem}

\clearpage

\begin{theorem}[$d_p$ Metric on $\mathbb{R}^2$]
    The metric space $\mathbb{R}^2$, called the \emph{Euclidean plane}, is obtained by
    taking the set of ordered pairs of real numbers,
    \[
    x=(\xi_1,\xi_2), \qquad y=(\eta_1,\eta_2),
    \]
    and defining the \emph{Euclidean metric} by
    \[
    d_p(x,y)=\left(|\xi_1-\eta_1|^p+|\xi_2-\eta_2|^p\right)^{\frac{1}{p}} \qquad (p\ge 1).
    \]
    Show that $d_p$ is a metric on $\mathbb{R}^2$.
\end{theorem}

\clearpage

\begin{theorem}[$d_{\infty}$ Metric on $\mathbb{R}^2$]
    The metric space $\mathbb{R}^2$, called the \emph{Euclidean plane}, is obtained by
    taking the set of ordered pairs of real numbers,
    \[
    x=(\xi_1,\xi_2), \qquad y=(\eta_1,\eta_2),
    \]
    and defining the \emph{Euclidean metric} by
    \[
    d_\infty(x,y)=\max\{|\xi_1-\eta_1|,|\xi_2-\eta_2|\}.
    \]
    Show that $d_{\infty}$ is a metric on $\mathbb{R}^2$.
\end{theorem}

\clearpage

\begin{heading}
    Pseudo-Metric on $\mathbb{R}^2$
\end{heading}

\vspace{10pt}

\begin{problem}[Pseudo-Metric on $\mathbb{R}^2$]
    The pseudo-metric on $\mathbb{R}^2$, is obtained by
    taking the set of ordered pairs of real numbers,
    \[
    x=(\xi_1,\xi_2), \qquad y=(\eta_1,\eta_2),
    \]
    and defining the \emph{Pseudo-Metric} by
    \[
    d(x,y)=|\xi_1-\eta_1|.
    \]
    Show that $d$ is a pseudo-metric on $\mathbb{R}^2$.
\end{problem}

\clearpage

\clearpage

\begin{heading}
    Metric Space $\mathbb{R}^n$
\end{heading}

\vspace{10pt}

\begin{theorem}[Standard Euclidean Metric in $\mathbb{R}^n$]
    The metric space $\mathbb{R}^n$ is obtained by taking the set of ordered n-tuples of
    real numbers,
    \[
    x=(\xi_1,\xi_2,\ldots,\xi_n), \qquad y=(\eta_1,\eta_2,\ldots,\eta_n),
    \]
    and defining the \emph{Euclidean metric} by
    \[
    d(x,y)=\left(\sum_{i=1}^{n}|\xi_i-\eta_i|^2\right)^{\frac{1}{2}} 
    \]
    Show that $d$ is a metric in $\mathbb{R}^n$.
\end{theorem}

\clearpage

\begin{theorem}[$d_p$ Metric in $\mathbb{R}^n$]
    The metric space $\mathbb{R}^n$ is obtained by taking the set of ordered n-tuples of
    real numbers,
    \[
    x=(\xi_1,\xi_2,\ldots,\xi_n), \qquad y=(\eta_1,\eta_2,\ldots,\eta_n),
    \]
    and defining the \emph{Euclidean metric} by
    \[
    d_p(x,y)=\left(\sum_{i=1}^{n}|\xi_i-\eta_i|^p\right)^{\frac{1}{p}} \qquad (p\ge 1).
    \]
    Show that $d_p$ is a metric in $\mathbb{R}^n$.
\end{theorem}

\clearpage

\begin{theorem}[$d_{\infty}$ Metric in $\mathbb{R}^n$]
    The metric space $\mathbb{R}^n$ is obtained by taking the set of ordered n-tuples of
    real numbers,
    \[
    x=(\xi_1,\xi_2,\ldots,\xi_n), \qquad y=(\eta_1,\eta_2,\ldots,\eta_n),
    \]
    and defining the \emph{Euclidean metric} by
    \[
    d_\infty(x,y)=\max\{|\xi_1-\eta_1|,|\xi_2-\eta_2|,\ldots,|\xi_n-\eta_n|\}.
    \]
    Show that $d_{\infty}$ is a metric in $\mathbb{R}^n$.
\end{theorem}

\clearpage

\begin{theorem}
Prove that the function $d_{\infty} : \mathbb{R}^n \times \mathbb{R}^n \to \mathbb{R}$
defined by
\[
d'(x,y)=\max_{1\le i\le n}|x_i-y_i|\]
where
\(x=(x_1,x_2,\dots,x_n),\;
y=(y_1,y_2,\dots,y_n)\in\mathbb{R}^n,
\)
is a metric on $\mathbb{R}^n$. Also prove that
\[
d_{\infty}(x,y)\le d(x,y)\le \sqrt{n}\,d_{\infty}(x,y),
\qquad \forall\, x,y\in\mathbb{R}^n,
\]
where $d$ denotes the standard Euclidean metric on $\mathbb{R}^n$.
\end{theorem}

\clearpage

\begin{heading}
    Metric Space $\mathbb{C}$
\end{heading}

\vspace{10pt}

\begin{theorem}[Standard metric on $\mathbb{C}$]
Prove that the function
\(
d : \mathbb{C} \times \mathbb{C} \to \mathbb{R}
\)
defined by
\begin{equation}
d(z_1,z_2)=|z_1-z_2|
\qquad \forall\, z_1,z_2\in\mathbb{C},
\end{equation}
forms a metric on $\mathbb{C}$.
\end{theorem}

\clearpage

\begin{theorem}[A metric on $\mathbb{C}$]
Prove that the function
\(
d : \mathbb{C} \times \mathbb{C} \to \mathbb{R}
\)
defined by
\begin{equation}
d(z_1,z_2)=\frac{|z_1-z_2|}{\sqrt{1+|z_1|^2}\sqrt{1+|z_2|^2}},
\qquad \forall\, z_1,z_2\in\mathbb{C},
\end{equation}
forms a metric on $\mathbb{C}$.
\end{theorem}

\clearpage

\begin{heading}
    Metric Space $\mathbb{C}^n$
\end{heading}

\vspace{10pt}

\begin{theorem}[Metric in $\mathbb{C}^n$]
Prove that the function
\(d : \mathbb{C}^n \times \mathbb{C}^n \to \mathbb{R}\)
defined by
\[
d(z,w)=\sqrt{\sum_{k=1}^{n}|z_k-w_k|^2}\]
where
\( z=(z_1,z_2,\dots,z_n),\;
w=(w_1,w_2,\dots,w_n)\in\mathbb{C}^n,
\)
forms a metric in $\mathbb{C}^n$.
\end{theorem}

\clearpage

\begin{heading}
    Function  Space $C[a,b]$
\end{heading}

\vspace{10pt}

\begin{theorem}[Max Metric on $C[a,b]$]
Prove that the function
\(d : C[a,b] \times C[a,b] \to \mathbb{R}\)
defined by
\[d(f,g)=\max_{x\in[a,b]}|f(x)-g(x)|\]
where   \(f,g\in C[a,b],\)
forms a metric on $C[a,b]$.
\end{theorem}

\clearpage

\begin{theorem}[Integral Metric on $C[a,b]$]
Prove that the function
\(d : C[a,b] \times C[a,b] \to \mathbb{R}\)
defined by
\[d(f,g)=\int_{a}^{b}|f(x)-g(x)|\,dx\]
where   \(f,g\in C[a,b],\)
forms a metric on $C[a,b]$.
\end{theorem}

\clearpage

\begin{heading}
    Space of Bounded Functions $B(A)$
\end{heading}

\vspace{10pt}

\begin{theorem}[Supremum Metric on $B(A)$]
Prove that the function
\(d : B(A) \times B(A) \to \mathbb{R}\)
defined by
\[d(f,g)=\sup_{x\in A}|f(x)-g(x)|\]
where   \(f,g\in B(A),\)
forms a metric on $B(A)$.
\end{theorem}

\clearpage

\begin{heading}
    Hamming Distance
\end{heading}

\vspace{10pt}

\begin{theorem}[Hamming Distance]
Prove that the function
\(d : X^n \times X^n \to \mathbb{R}\)
defined by
\[d(x,y)=\sum_{i=1}^{n}\delta(x_i,y_i)\]
where
\[\delta(a,b)=\begin{cases}
0, & \text{if } a=b;\\
1, & \text{if } a\ne b,
\end{cases}\]
and \(x=(x_1,x_2,\dots,x_n),\;
y=(y_1,y_2,\dots,y_n)\in X^n,\)
forms a metric on $X^n$.
\end{theorem}





\EndPage
\end{document}